% When submitting your files, remember to upload this *tex file, the pdf generated with it, the *bib file (if bibliography is not within the *tex) and all the figures.

% NB logo1.jpg is required in the path in order to correctly compile front page header %%%

\documentclass[utf8]{FrontiersinHarvard} % for articles in journals using the Harvard Referencing Style (Author-Date), for Frontiers Reference Styles by Journal: https://zendesk.frontiersin.org/hc/en-us/articles/360017860337-Frontiers-Reference-Styles-by-Journal
\usepackage{url,hyperref,lineno,microtype,subcaption}
%\usepackage[]{natbib, booktabs, graphicx, tabularx, makecell}
\usepackage{lscape, float, multirow, adjustbox}
\usepackage[T1]{fontenc}
\usepackage[onehalfspacing]{setspace}

\linenumbers


% Leave a blank line between paragraphs instead of using \\


\def\keyFont{\fontsize{8}{11}\helveticabold }
\def\firstAuthorLast{Morris {et~al.}} %use et al only if is more than 1 author
\def\Authors{Joanna Morris\,$^{1,2,*}$, Emma Kealey\,$^{2,3}$, Frankie Greene\,$^{1,4}$, Joemari Pulido\,$^{2,5}$, Tess Rooney\,$^{2,6}$, Ethan Moore\, $^{2}$, Crismar Ramos-Marte\, $^{2}$, Natalia Alzate\,$^{2}$}
% Affiliations should be keyed to the author's name with superscript numbers and be listed as follows: Laboratory, Institute, Department, Organization, City, State abbreviation (USA, Canada, Australia), and Country (without detailed address information such as city zip codes or street names).
% If one of the authors has a change of address, list the new address below the correspondence details using a superscript symbol and use the same symbol to indicate the author in the author list.
\def\Address{$^{1}$School of Cognitive Science, Hampshire College, Amherst, MA, USA \\
	$^{2}$Department of Psychology, Providence College, Providence, RI, USA \\
	$^{3}$Department of Linguistics, University of Southern California, Los Angeles, CA, USA \\
	$^{4}$Oregon Health \& Science University and Portland State University School of Public Health, Portland, OR, USA \\
	$^{5}$Department of Cognitive, Linguistic \& Psychological Sciences, Brown University, Providence, RI, USA \\
	$^{6}$Columbia University School of Nursing, New York, NY, USA}
% The Corresponding Author should be marked with an asterisk
% Provide the exact contact address (this time including street name and city zip code) and email of the corresponding author
\def\corrAuthor{Joanna Morris}

\def\corrEmail{jmorris6@providence.edu}


\begin{document}
	\onecolumn
	\firstpage{1}
	
	\title[Running Title]{Neural signatures of visual statistical learning dissociate from explicit familiarity judgments} 
	
	\author[\firstAuthorLast ]{\Authors} %This field will be automatically populated
	\address{} %This field will be automatically populated
	\correspondance{} %This field will be automatically populated
	
	\extraAuth{}% If there are more than 1 corresponding author, comment this line and uncomment the next one.
	%\extraAuth{corresponding Author2 \\ Laboratory X2, Institute X2, Department X2, Organization X2, Street X2, City X2 , State XX2 (only USA, Canada and Australia), Zip Code2, X2 Country X2, email2@uni2.edu}
	
	
	\maketitle
	



\begin{abstract}

%%% Leave the Abstract empty if your article does not require one, please see the Summary Table for full details.
\section{}
	Skilled reading requires us to rapidly retrieve a word's meaning from memory based on 	Skilled reading relies on sensitivity to morphological structure and the quality of lexical representations, with readers differing in how efficiently morphological information is processed over time. To examine how these differences unfold during word recognition, event-related potentials (ERPs) and reaction times were recorded while participants performed a lexical decision task on words and morphologically complex nonwords. Morphological structure was manipulated along multiple dimensions, including base frequency, morphological family size, and morphological complexity. Individual differences were quantified using continuous measures derived from a principal components analysis of vocabulary knowledge, spelling ability, and print exposure, yielding two orthogonal dimensions: Language Proficiency and Orthographic Sensitivity. Behaviorally, response times were primarily driven by item-level lexical familiarity and morphological complexity, with little evidence that individual differences selectively modulated morphological effects. In contrast, ERP measures revealed robust, stage-specific modulation by reader skill. In contrast, ERP measures revealed clear, stage-specific effects. The N250 showed strong modulation by individual differences during processing of morphologically complex items, whereas N400 effects were selectively shaped by Language Proficiency when morphological structure  present in the absence of a lexical entry. These findings indicate that individual differences in reading skill are expressed more clearly in neural processing over time than in overt behavior, highlighting the importance of electrophysiological measures for understanding how morphology is parsed and regulated during skilled reading.


\tiny
 \keyFont{ \section{Keywords:} morphological processing, individual differences, event-related potentials, N250, N400, lexical decision} %All article types: you may provide up to 8 keywords; at least 5 are mandatory.
\end{abstract}

\section{Introduction}

Reading is often described as “the process of gaining meaning from print” \citep[][p.~34]{raynerHowPsychologicalScience2001}. Although comprehension ultimately depends on sentence- and discourse-level integration, efficient access to individual word forms and meanings remains a foundational component of skilled reading. A central challenge for readers is therefore to extract structure from written words rapidly and flexibly, using information distributed across orthography, morphology, and semantics.

In languages such as English, this challenge is amplified by the properties of the writing system itself. English exhibits irregular and context-sensitive grapheme–phoneme mappings, rendering phonological decoding alone an insufficient basis for fluent word recognition. At the same time, English preserves regularities at higher representational levels, particularly in its morphological structure. Words that share a stem or affix often exhibit consistent orthographic patterns and systematic semantic relationships (e.g., heal–health, teach–teacher). As a result, skilled reading in English requires sensitivity not only to phonological correspondences but also to recurring morphological structure that links form and meaning.

The question of how morphological structure is represented and processed during word recognition has generated substantial empirical and theoretical debate. A core issue concerns the nature of the units involved: whether skilled readers decompose complex words into their constituent morphemes during recognition, access whole-word representations directly, or exploit learned statistical regularities that render such architectural distinctions unnecessary. These possibilities are not merely theoretical alternatives—they carry different implications for what it means to be a skilled morphological processor, and for how individual differences in reading ability should modulate the recognition of morphologically complex words. Adjudicating among them requires identifying variables that are selectively diagnostic of the proposed mechanisms, and two such variables—stem frequency and morphological family size—have proven particularly central to this debate.

\subsection{Stem Frequency and Morphological Family Size: Theoretical Interpretations Across Current Models}

Stem frequency refers to the summed token frequency of a base morpheme across all its surface realizations—inflected and derived forms alike—and is often operationalized as lemma frequency or cumulative root frequency [CITATION]. Morphological family size refers to the number of distinct word types that share a base morpheme, encompassing both compounds and derivatives [CITATION]. Crucially, the two variables are empirically dissociable: a base can be high in stem frequency but low in family size, or vice versa. Both facilitate visual lexical decision and naming, but their effects are independent when appropriately controlled [CITATION], and it is this dissociation that has driven theoretical debate. What each effect represents depends substantially on which processing architecture one adopts, and current models differ not only in their answers but in the kinds of representations they permit as answers.

In the classical framework developed by Baayen and Schreuder [CITATION] and related dual-route accounts, the two effects were treated as converging diagnostics of separate processing routes. Stem frequency was held to reflect the ease of accessing a stored morpheme entry via a decompositional route: because complex words are parsed into their constituent morphemes during recognition, the frequency of the stem as an independent lexical unit facilitates that parsing step and any subsequent form built upon it. Morphological family size, by contrast, was attributed to the whole-form route and to the richness of the semantic network surrounding a base morpheme: encountering a word activates not only its own representation but those of its morphological relatives, and a larger family produces broader co-activation of the base representation [CITATION]. On this view, the dissociation between the two effects was itself evidence for a dual-route architecture, with each variable selectively loading on one route.

Subsequent models have substantially revised these interpretations, though they differ in the depth of revision they undertake. The early automatic decomposition account developed by Rastle, Davis, and colleagues [CITATION], drawing on masked morphological priming, proposes that decomposition is not one route among several but an obligatory initial parsing step that applies to all morphologically structured—and apparently morphologically structured—inputs regardless of semantic transparency. On this account, stem frequency reflects the ease of this early parsing stage: higher stem frequency facilitates the activation of the morpheme representation that is the immediate output of decomposition. Morphological family size is then attributed to a later, semantically sensitive stage, consistent with the finding that opaque relatives contribute less to the family size effect than transparent ones [CITATION]. The two-stage sequential architecture thus preserves the dissociation between the two effects but reconceptualizes it: rather than reflecting competition between parallel routes, the effects index distinct temporal stages of a serial process, with stem frequency effects arising early and obligatorily and family size effects arising later and conditionally on semantic coherence.

Grainger and colleagues' morpho-orthographic account [CITATION] offers a partially overlapping but architecturally distinct interpretation. On this view, an intermediate level of representation between sublexical letter clusters and whole-word form representations is active in skilled reading, at which morphological constituents are represented as units—not as stored symbolic morphemes, but as learned orthographic patterns that reliably co-occur with semantic content. Stem frequency on this account indexes the strength of representation at this morpho-orthographic level: a high-frequency stem has a more robustly activated intermediate representation, facilitating subsequent access to meaning. Morphological family size, in turn, is reflected in the richness of the semantic representation that a morpho-orthographic unit connects to—a larger family means a more densely structured semantic neighborhood, supporting faster and more confident lexical retrieval. This account differs from the early decomposition account in treating the intermediate representations as learned, gradient, and form-based rather than as the output of a parsing procedure, but like the early decomposition account it assigns the two effects to different levels of a multi-stage architecture.

The most radical reinterpretation is offered by the discriminative learning framework, instantiated in Naive Discriminative Learning [NDL; CITATION] and its successor, Linear Discriminative Learning [LDL; CITATION]. These models share the commitment that morphological processing involves no stored morpheme representations, no decomposition procedure, and no dedicated morphological architecture of any kind. Instead, the system learns a direct mapping from distributional form cues—typically character n-grams—to semantic representations, via error-driven weight adjustment. Within this framework, stem frequency and family size effects are reframed not as diagnostics of separate routes or stages, but as emergent consequences of a single learning mechanism operating over different aspects of distributional structure.

Stem frequency, in NDL and LDL, reflects the accumulated weight strength of the n-gram cues that spell out the base morpheme. Because these cues appear across all surface forms of the base—across every inflected and derived word in which the stem is present—they accumulate training signal proportional to the base's cumulative token frequency. High stem frequency therefore means that the relevant form cues have been more extensively and precisely calibrated, producing stronger and more reliable activation of the associated semantic representation. Critically, this occurs without any discrete stem entry being accessed: the effect is a gradient property of the learned weight matrix rather than the consequence of retrieving a stored morpheme. Morphological family size, by contrast, reflects a different statistical property of the same learning environment: the type diversity of the training signal rather than its token accumulation. A large morphological family exposes the stem's n-gram cues to a wide variety of distinct form-meaning pairs, each sharing the same semantic core but differing in their idiosyncratic features. The error-driven learning dynamics of NDL—and the analogous optimization pressure in LDL's matrix learning—cause the weights for shared cues to converge on the semantic dimensions that are consistent across the family, while suppressing idiosyncratic variation. A large, semantically coherent family therefore produces a more robustly and precisely represented semantic core, facilitating lexical retrieval. The well-documented finding that semantically opaque relatives fail to contribute to the family size effect [CITATION] follows without additional stipulation: an opaque relative has a semantic vector that is distant from the base in distributional space, and therefore provides no consistent reinforcement of the stem's core semantic weights during learning.

The theoretical significance of the NDL/LDL account lies not only in its reframing of each effect individually, but in its claim that both effects arise from a single mechanism. Where dual-route and multi-stage models took the dissociation between stem frequency and family size as evidence for distinct representational levels or processing routes, the discriminative learning framework attributes the dissociation to two aspects of the same distributional statistics—token frequency and type diversity—acting through the same learned weight structure. This is a substantively more parsimonious account, though it comes at the cost of reconceptualizing what the effects are evidence for: rather than revealing the architecture of a morphological processor, they reveal properties of the statistical environment in which the lexicon was learned.

These competing interpretations carry importantly different predictions about the time course of stem frequency and family size effects, about how the system responds to novel morphological structure, and about the role of individual differences in processing efficiency—three dimensions along which the present study is specifically designed to adjudicate.

\subsection{Individual Differences in Morphological Processing}

A complementary perspective on the architecture of morphological processing comes from examining variation across readers. Traditional models of word recognition have typically characterised processing at the level of the skilled reader in general, abstracting over individual differences in favour of a uniform cognitive architecture. Yet if morphological sensitivity reflects the accumulated product of statistical learning rather than the output of a fixed parsing mechanism, readers who differ in the depth or precision of their lexical representations should differ systematically in how they process morphologically complex words—and potentially in which aspects of that processing are most affected.

The Lexical Quality Hypothesis [CITATION] provides a principled account of such differences. On this view, lexical representations vary in quality across readers—that is, in how completely and precisely a word's orthographic, phonological, and semantic constituents are specified and interconnected. High-quality representations enable rapid, automatic activation of a word's sound and meaning from its visual form; low-quality representations are partial or underspecified, producing processing that is slower, less reliable, and more dependent on contextual support. Individual differences in reading and spelling skill can be understood as reflecting differences in lexical quality [CITATION], with more skilled readers possessing more precise orthographic representations that support faster word identification and greater resistance to interference from orthographically similar forms.

Critically, readers do not differ along a single dimension of lexical quality. Andrews and colleagues [CITATION] have demonstrated that individual differences in orthographic and semantic knowledge are dissociable and predict distinct aspects of lexical processing. Readers whose spelling ability exceeds their vocabulary knowledge—an orthographic profile reflecting high precision in sublexical form representations—show robust morphological priming even for semantically opaque pairs such as corner–corn, suggesting that their processing is driven primarily by form-based segmentation at the level of morpho-orthographic structure. Readers whose vocabulary knowledge exceeds their spelling ability—a semantic profile reflecting richer lexical-semantic representations—show morphological priming only for transparent pairs, indicating that their processing is more strongly constrained by meaning. These dissociable profiles suggest that the relative contributions of form-based and meaning-based information to morphological processing are not fixed properties of the recognition system but vary systematically across readers as a function of the type of lexical knowledge they have accumulated.

This dissociation maps directly onto the theoretical contrast between the processing accounts reviewed above. Readers with a strong orthographic profile resemble the system described by morpho-orthographic and decompositional accounts—their processing is driven by sensitivity to sublexical form structure, the primary cue in discriminative learning models and the basis for early automatic decomposition in parsing accounts. Readers with a strong semantic profile resemble the system in which family size and semantic coherence are the dominant influences, consistent with the later semantic stage proposed by multi-stage accounts and with the meaning-space organisation of LDL. Rather than treating this as an either-or contrast between reader types, the present study operationalises these profiles as continuous dimensions derived from a principal components analysis of vocabulary knowledge, spelling ability, and print exposure, yielding two orthogonal components—Language Proficiency and Orthographic Sensitivity—that capture the full range of the dissociation in the sample. If the effects of stem frequency and family size are modulated differently by these two dimensions, and if those modulations emerge at distinct temporal stages in the ERP, this would constitute strong evidence that the theoretical distinction between form-based and meaning-based morphological processing is not merely a property of competing computational models but a genuine axis of variation in the human reading system.

\subsection{Aims and Predictions}

Despite robust evidence that morphological structure shapes visual word recognition, the theoretical debate between structured parsing accounts and discriminative learning frameworks has remained difficult to resolve at the level of group-averaged data. Both classes of model can accommodate the core behavioral effects—facilitation by stem frequency and morphological family size—and both can be reconciled with the broad temporal dissociation between early form-based and later semantic processing. What they cannot equally accommodate are fine-grained differences in when and for whom these effects emerge, and whether the system generalizes morphological regularities productively to novel forms. Our primary aim is therefore to use individual differences as a diagnostic tool: if morphological sensitivity reflects accumulated statistical learning rather than fixed representational primitives, then Orthographic Sensitivity and Language Proficiency should modulate morphological effects in ways that are temporally dissociable—with Orthographic Sensitivity predicting early form-level processing and Language Proficiency predicting later semantic processing—rather than scaling uniformly with general language experience. This predicted dissociation constitutes the study's central theoretical claim, and it is evaluated through two complementary analyses.

The first analysis examines the time course of stem frequency and family size effects in real word recognition. Because the real words in the present study were not matched across morphological variables—reflecting the naturalistic distributional properties of morphologically complex words rather than an orthogonally controlled factorial design—this analysis treats stem frequency, family size, and morphological complexity as continuous predictors in a regression framework for reaction time, examining their independent contributions across the full item set. For ERP data, signal averaging requirements necessitate a different approach: items were median-split into high and low groups separately for stem frequency and family size, and ERPs were averaged within each group to yield stable waveforms for comparison. The two dependent measures therefore carry somewhat different inferential weight: the RT regression speaks to the independent continuous contributions of each variable across the full distributional range, while the ERP analyses speak to the direction and timing of effects at the extremes of each variable's distribution. Multi-stage architectures—whether decompositional [CITATION] or morpho-orthographic [CITATION]—predict that stem frequency and family size should influence processing at distinct temporal stages, with stem frequency modulating earlier components reflecting form-level access and family size exerting its primary influence later when semantic information becomes available. The discriminative learning framework [CITATION], by contrast, treats both effects as arising from a single learned weight structure, predicting no systematic difference in their temporal onset or scalp distribution. We predicted that this temporal dissociation, if present, would be most pronounced in readers high in Language Proficiency—for whom family size should carry particular weight—and that early stem frequency effects would be most strongly modulated by Orthographic Sensitivity.

The second analysis examines generalization to novel morphological structure using morphologically complex nonwords. Because morphologically complex nonwords have no stored whole-word representation, they place maximal pressure on whatever generalization mechanism the system possesses and allow stem-level variables to be examined independently of whole-word familiarity—a control that the real word analysis cannot straightforwardly provide. Morphologically structured nonwords (gasful: real stem + real affix) were compared against matched controls (gasfil: real stem + pseudoaffix), holding stem identity constant while manipulating the morphological legitimacy of the suffix. Rule-governed parsing accounts predict that this manipulation should yield early ERP differences, specifically modulation of the N250, reflecting the detection of constituent structure at the suffix boundary—an effect that should be present categorically wherever a valid parse is licensed and should not vary systematically across readers of comparable literacy. The discriminative learning framework also predicts sensitivity to the suffix manipulation, since real affixes carry stronger and more consistent n-gram cue mappings to morphological outcomes than pseudoaffixes do, but attributes this to the same learned weight structure underlying real word recognition rather than to a dedicated parsing operation. The present design is therefore better positioned to evaluate where in time and for whom suffix legitimacy influences processing than to adjudicate between categorical and graded generalization directly. Critically, if Orthographic Sensitivity modulates the magnitude or onset of the N250 response to the suffix contrast, this would implicate form-cue learning rather than a fixed parsing operation, since the output of a rule-governed parser should be invariant across readers who have all clearly acquired basic morphological knowledge.

Taken together, the two analyses ask whether the distinction between form-based and meaning-based morphological processing—long debated as a property of competing computational architectures—is also a genuine axis of variation in the human reading system. If Orthographic Sensitivity selectively modulates early ERP effects of stem frequency and suffix structure while Language Proficiency selectively modulates later N400 effects of family size, the data would support a view in which different types of accumulated lexical knowledge contribute to qualitatively distinct stages of morphological recognition, with consequences for how individual differences in reading should be theorised and measured.

\section{Methods}

\subsection{Materials}

The experimental stimulus set included 200 nonwords, constructed to manipulate morphological structure while controlling surface form. In the \textit{complex nonword} condition, we combined existing English base words (e.g., \textit{gas}) with existing derivational suffixes (e.g., \textit{gasful}), to yield 100 morphologically plausible complex nonwords. Each combination was syntactically legal and reflected common English derivational patterns. A total of 16 suffixes were used, each attached to four different stems.

In the \textit{simple nonword} condition, we cobmined thesame base words as in the complex non-word condition  with non-suffix endings that were orthographically similar to the real suffixes used in the complex condition but did not constitute valid morphemes, to yield 100 orthographically plausible simple non-words. These non-affix endings were created by changing the central letter of each suffix (e.g., \textit{gasful}–\textit{gasfil}), preserving length and orthographic neighborhood characteristics while eliminating morphological structure.  Thus each complex non-word had an orthographically matched simple counterpart.

Because the same base words appeared in both the simple and complex conditions, the experimental nonwords were distributed across two lists such that no participant encountered the same base more than once. In addition to the experimental nonwords, the stimulus set included 100 monomorphemic fillers (50 real words and 50 nonwords) to balance lexical status and reduce predictability.

Estimates of cumulative base (root) frequency, surface frequency, and morphological family size were obtained from the CELEX lexical database. Morphological family size was calculated by identifying all morphologically related forms sharing the same base (e.g., \textit{calculate}, \textit{calculable}, \textit{calculation}, \textit{calculator}). Following prior work \citep[e.g.,][]{deJong2001, deJong2003}, both compounds (e.g., \textit{WATCHTOWER}) and hyphenated compounds (e.g., \textit{CHECK-IN}) were included, as they contribute to morphological connectivity within the lexicon. Base frequency was computed by summing the lemma frequencies of all confirmed family members.

For reaction time analyses, all lexical and individual-difference variables were treated as continuous predictors to preserve information and maximize statistical power. In contrast, ERP analyses were conducted on condition-averaged waveforms to improve signal-to-noise ratio; as a result, item-level predictors such as base frequency and family size were operationalized as categorical condition factors rather than continuous trial-level covariates. This distinction reflects standard practice in ERP research and aligns the statistical treatment of predictors with the structure of the underlying data.

\subsection{Individual Differences and Grouping Variables}

To quantify individual differences in linguistic experience and orthographic knowledge, participants completed three standardized measures: the Nelson–Denny Vocabulary Test \citep{nelson1960}, a spelling task adapted from \citet{sawi2016}, and the Author Recognition Test (ART; \citealt{stanovich1989}). Scores on these measures were moderately correlated ($r$s = .29–.51, all $p < .02$), indicating partially overlapping but dissociable components of lexical experience.

To capture the shared and distinct variance among these measures, we conducted a principal components analysis (PCA) on standardized (z-scored) vocabulary, spelling, and ART scores. PCA was selected to derive orthogonal dimensions of individual differences without imposing a priori weighting or categorical group boundaries. Two components had eigenvalues greater than 1, jointly accounting for 83.5\% of the total variance (58.0\% for Dimension~1 and 25.5\% for Dimension~2).

Component loadings (Table~\ref{tab:pca_loadings}) indicated that Dimension~1 loaded most strongly on vocabulary and print exposure (ART), with a secondary contribution from spelling accuracy, and was interpreted as indexing \textit{Language Proficiency}. Dimension~2 loaded most strongly on spelling accuracy, with weaker and oppositely signed contributions from vocabulary and ART, and was interpreted as indexing \textit{Orthographic Sensitivity}:

\begin{itemize}
	\item[-] \textit{Dimension~1 (Language Proficiency)}: vocabulary (.82), ART (.81), spelling (.65)
	\item[-] \textit{Dimension~2 (Orthographic Sensitivity)}: spelling (.76), vocabulary ($- .29$), ART ($- .32$)
\end{itemize}

The two components were statistically independent ($r \approx 0$), supporting their interpretation as orthogonal dimensions of reading-related skill. These continuous component scores were used as predictors in all behavioral and ERP analyses.


\begin{table}[ht]
	\centering
	\caption{PCA loadings and variance explained for the three lexical-experience measures.}
	\label{tab:pca_loadings}
	\begin{tabular}{lcc}
		\toprule
		\textbf{Measure} & \textbf{Dimension~1 (Language Proficiency)} & \textbf{Dimension~2 (Orthographic Sensitivity)} \\
		\midrule
		Vocabulary (z\_vcb) & 0.82 & $-0.29$ \\
		Spelling (z\_spl) & 0.65 & 0.76 \\
		Author Recognition Test (z\_art) & 0.81 & $-0.32$ \\
		\midrule
		\textit{Eigenvalue} & 1.74 & 0.76 \\
		\textit{Variance explained (\%)} & 58.0 & 25.5 \\
		\bottomrule
	\end{tabular}
\end{table}

\section{Statistical Analyses}

The analyses are organized around the study’s primary theoretical aim: determining whether Orthographic Sensitivity and Language Proficiency dissociably modulate the timing and nature of morphological processing. We predicted that Orthographic Sensitivity would selectively modulate early, form-based processing indexed by the N250, while Language Proficiency would selectively modulate later semantic integration indexed by the N400. The behavioral and ERP results are reported separately for words and nonwords in each section, followed by a summary that evaluates these predictions against the obtained pattern.

\subsection{Behavioral Data (Reaction Times)}
Reaction time (RT) analyses were conducted on correct responses only. RTs were log-transformed to reduce positive skew and to better meet model assumptions of normality and homoscedasticity. Separate analyses were conducted for word and nonword trials, reflecting differences in the availability and interpretation of morphological predictors across lexicality.

RTs were analyzed using linear mixed-effects models fit with the lme4 package in R. All models included random intercepts for participants and items, and random slopes for within-participant predictors where supported by the data. Continuous lexical predictors (log base frequency and log family size) were mean-centered and scaled prior to analysis. Individual-difference measures—Orthographic Sensitivity and Language Proficiency—were entered as continuous predictors corresponding to standardized PCA scores.

For word trials, models tested the effects of base frequency, family size, Orthographic Sensitivity, and Language Proficiency, as well as interactions between orthographic sensitivity and lexical familiarity measures. For nonword trials, models additionally included morphological complexity (simple vs. complex) as a categorical predictor and tested whether complexity effects were modulated by individual differences. Treatment (reference) coding was used for categorical predictors in RT analyses. Model significance was assessed using Type III tests with Satterthwaite-approximated degrees of freedom. Fixed-effect estimates are reported alongside model-based estimated marginal means where relevant.

\subsection{ERP Data}
ERP analyses were conducted on mean amplitudes computed from condition-averaged waveforms. Separate datasets were constructed for words and nonwords, and analyses were conducted independently for the N250 and N400 time windows. Because ERP amplitudes were derived from condition-level averages rather than trial-level data, morphological predictors in ERP analyses were treated as categorical condition factors rather than continuous covariates.

ERP amplitudes were analyzed using linear mixed-effects models with fixed effects of morphological structure (family size for words; complexity and family size for nonwords), Orthographic Sensitivity, and Language Proficiency, as well as their interactions. Models included random intercepts for participants and for electrode site nested within participant, allowing baseline amplitude differences across electrodes to be modeled while preserving the spatial structure of the data.

Categorical predictors in ERP analyses were effect-coded (sum-to-zero coding), such that model intercepts represent grand mean amplitudes and main effects reflect average effects across conditions. Orthographic Sensitivity and Language Proficiency were entered as continuous, mean-centered predictors. Time window was not included as a factor in the models; instead, separate models were fit for the N250 and N400 to preserve component-specific interpretation.

Statistical significance of fixed effects was assessed using Type III F-tests with Satterthwaite-approximated degrees of freedom. Where higher-order interactions involving individual-difference measures were significant, follow-up analyses were conducted using estimated marginal means (emmeans). In particular, interactions involving Language Proficiency were probed by estimating contrasts at ±1 SD of the proficiency distribution, allowing direct tests of whether morphological effects differed as a function of proficiency.

\subsection{Model Estimation and Inference}
All mixed-effects models were fit using restricted maximum likelihood estimation. To ensure stable convergence, models were estimated using the BOBYQA optimizer. Continuous predictors were centered to facilitate interpretation of interaction terms and model intercepts. For predictors with one degree of freedom, F-tests from the analysis-of-variance tables are mathematically equivalent to squared t-tests from model summaries; all inferential conclusions are therefore invariant to the choice of reporting format.

Across all analyses, inferential emphasis was placed on theoretically motivated interactions rather than exhaustive post hoc comparisons. Follow-up contrasts were conducted only where warranted by significant interactions and were limited to tests directly addressing the study’s hypotheses.

\section{Results}

The analyses are organized around the study's primary theoretical aim: determining whether Orthographic Sensitivity and Language Proficiency dissociably modulate the timing and nature of morphological processing. We predicted that Orthographic Sensitivity would selectively modulate early, form-based processing indexed by the N250, while Language Proficiency would selectively modulate later semantic integration indexed by the N400. The behavioral and ERP results are reported separately for words and nonwords in each section, followed by a summary that evaluates these predictions against the obtained pattern.

\subsection{Words}

\paragraph*{Reaction Times }
Reaction times (RTs) for correct lexical decisions to word stimuli were analyzed using a linear mixed-effects model with fixed effects of Base Frequency, Family Size, Orthographic Sensitivity (OS), and Language Proficiency (LP), as well as interactions between OS and Base Frequency and Family Size. The model included random intercepts for participants and items, as well as by-participant random slopes for Base Frequency and Family Size. The analysis revealed robust main effects of item-level lexical familiarity (Base Frequency and Family Size). Words derived from higher-frequency bases were recognized more quickly than those derived from lower-frequency bases, $\beta = -0.018$, $SE = 0.005$, $t(101.85) = -3.69$, $p < .001$. Similarly, words from larger morphological families elicited faster responses than those from smaller families, $\beta = -0.013$, $SE = 0.005$, $t(93.20) = -2.60$, $p = .011$. These effects were independent: neither the Base Frequency × Orthographic Sensitivity interaction, $\beta = -0.002$, $SE = 0.003$, $t(191.85) = -0.93$, $p = .355$, nor the Family Size × Orthographic Sensitivity interaction, $\beta = 0.001$, $SE = 0.003$, $t(78.22) = 0.44$, $p = .66$, approached significance. Orthographic Sensitivity showed a marginal trend toward faster overall response times, $\beta = -0.035$, $SE = 0.019$, $t(63.88) = -1.85$, $p = .069$, whereas Language Proficiency did not exert a main effect on RTs, $\beta = 0.013$, $SE = 0.012$, $t(63.97) = 1.05$, $p = .30$. Importantly, neither individual-difference measure selectively modulated the effects of Base Frequecy or Family Size on response speed. Together, these results indicate that response times to words are driven primarily by iBase Frequency and Family Size, with only weak evidence that individual differences influence overall efficiency. There was no evidence that Orthographic Sensitivity or Language Proficiency alters the structure of morphological effects on word recognition speed.

\paragraph*{Words (N250)}
Mean N250 amplitudes for word stimuli were analyzed using a linear mixed-effects model with fixed effects of Family Size, Orthographic Sensitivity (OS), and Language Proficiency (LP), including interactions between Family Size and each individual-difference measure. The model included random intercepts for participants and for electrode site nested within participant. The analysis revealed a robust main effect of Family Size, such that words from larger morphological families elicited more positive N250 amplitudes than words from smaller families, $\beta = 0.160$, $SE = 0.029$, $t(1617) = 5.48$, $p < .001$. Neither Orthographic Sensitivity, $\beta = 0.261$, $SE = 0.256$, $t(57) = 1.02$, $p = .31$, nor Language Proficiency, $\beta = -0.279$, $SE = 0.186$, $t(57) = -1.50$, $p = .14$, yielded significant main effects. Interactions between Family Size and individual-difference measures were not reliable. Neither the Family Size × Orthographic Sensitivity interaction, $\beta = 0.055$, $SE = 0.033$, $t(1617) = 1.67$, $p = .095$, nor the Family Size × Language Proficiency interaction were significant, $\beta = -0.036$, $SE = 0.024$, $t(1617) = -1.53$, $p = .13$. Together, these results indicate that morphological family structure modulates early morpho-orthographic processing for words at the N250, with limited evidence that individual differences shape this effect at this processing stage. Notably, the effect of Family Size at the N250 — an early component typically associated with form-level processing — is not straightforwardly predicted by multi-stage accounts, which attribute family size effects primarily to later semantic stages. This finding may indicate that morphological family structure influences the efficiency of early orthographic processing for familiar words, possibly reflecting the contribution of family-level distributional statistics to the strength of morpho-orthographic representations, though this interpretation requires further investigation.

\subsection*{Words N400}
Mean N400 amplitudes for word stimuli were analyzed using a linear mixed-effects model with fixed effects of Family Size, Orthographic Sensitivity (OS), and Language Proficiency (LP), including interactions between Family Size and each individual-difference measure. Random intercepts were included for participants and for electrode site nested within participant. The analysis revealed a robust main effect of Family Size, such that words from larger morphological families elicited less negative N400 amplitudes than words from smaller families, $\beta = 0.224$, $SE = 0.032$, $t(1617) = 7.01$, $p < .001$. This effect indicates facilitated semantic integration for words supported by richer morphological families. In addition, a significant main effect of Orthographic Sensitivity was observed, $\beta = 0.677$, $SE = 0.318$, $t(57) = 2.13$, $p = .038$, reflecting overall differences in N400 amplitude as a function of spelling-based reading skill. This main effect of Orthographic Sensitivity — reflecting overall differences in N400 amplitude rather than selective modulation of family size effects — was not predicted and suggests that spelling-based reading skill influences the general efficiency of semantic access for familiar words independently of the specific morphological properties of individual items. In contrast, Language Proficiency did not yield a main effect, $\beta = 0.001$, $SE = 0.231$, $t(57) = 0.01$, $p = .996$. Critically, neither the Family Size × Orthographic Sensitivity interaction, $\beta = 0.044$, $SE = 0.036$, $t(1617) = 1.24$, $p = .214$, nor the Family Size × Language Proficiency interaction, $\beta = -0.035$, $SE = 0.026$, $t(1617) = -1.34$, $p = .180$, reached significance. Thus, while both morphological family structure and orthographic sensitivity independently influenced N400 amplitudes, there was limited evidence that individual differences modulated the magnitude of family-size effects at this later stage of processing for familiar words.



\subsection{ERP Results: Neural Indices of Morphological Processing}
While reaction times provide a sensitive measure of overall processing efficiency, they do not reveal how morphological information is represented and integrated over time. To examine whether individual differences modulate the temporal dynamics of morphological processing—beyond what is evident in behavioral responses—we next analyzed event-related potentials (ERPs), focusing on the N250 and N400 components as indices of early morpho-orthographic parsing and later semantic integration, respectively. Separate models were fit for words and nonwords to account for differences in lexical status and available morphological cues. We tested whether morphological structure interacts with Orthographic Sensitivity and Language Proficiency at each processing stage, allowing individual differences to influence early parsing, later integration, or both.

\subsubsection{Nonwords}

\paragraph*{Reaction Times}
Reaction times (RTs) for correct lexical decisions to nonwords were analyzed using a linear mixed-effects model with fixed effects of Morphological Complexity (simple vs.\ complex), Base Frequency, Family Size, Orthographic Sensitivity (OS), and Language Proficiency (LP), as well as the theoretically motivated Complexity × Orthographic Sensitivity interaction. The model included random intercepts for participants and items and by-participant random slopes for Complexity. The analysis revealed a strong main effect of Morphological Complexity: complex nonwords elicited significantly slower responses than simple nonwords, $\beta = -0.049$, $SE = 0.005$, $t(61.58) = -9.09$, $p < .001$. In addition, Base Frequency exerted a reliable inhibitory effect, such that nonwords derived from higher-frequency bases were responded to more slowly than those from lower-frequency bases — a pattern consistent with discriminative learning accounts, in which high-frequency stems generate strong sublexical cue activation that competes with, and thereby delays, the rejection decision for nonword forms. $\beta = 0.015$, $SE = 0.005$, $t(98.23) = 3.05$, $p = .003$. Family Size did not significantly influence response times, $\beta = -0.006$, $SE = 0.005$, $t(96.65) = -1.24$, $p = .22$. Orthographic Sensitivity showed a significant main effect, with higher-sensitivity participants responding more quickly overall, $\beta = -0.042$, $SE = 0.019$, $t(61.52) = -2.25$, $p = .028$. In contrast, Language Proficiency did not yield a main effect on nonword RTs, $\beta = 0.003$, $SE = 0.012$, $t(62.91) = 0.27$, $p = .79$. Critically, Orthographic Sensitivity did not interact with Morphological Complexity, $\beta = 0.002$, $SE = 0.006$, $t(59.95) = 0.37$, $p = .71$, indicating that greater orthographic skill did not selectively reduce the processing cost associated with morphologically complex nonwords. Taken together, these results indicate that nonword decision latencies are strongly shaped by morphological complexity and item-level lexical familiarity, particularly base frequency, while individual differences primarily affect overall response efficiency. There was no evidence that orthographic or language proficiency alters the magnitude of morphological complexity effects in behavioral responses to novel forms.

\paragraph*{Nonwords N250}
Mean N250 amplitudes for nonwords were analyzed using a linear mixed-effects model with fixed effects of Morphological Complexity (simple vs.\ complex), Family Size, Orthographic Sensitivity (OS), and Language Proficiency (LP), including all theoretically motivated interactions. Random intercepts were included for participants and for electrode site nested within participant. The analysis revealed significant main effects of both Morphological Complexity and Family Size. Complex nonwords elicited larger N250 amplitudes than simple nonwords, $\beta = 0.116$, $SE = 0.025$, $t(4851) = 4.64$, $p < .001$, and nonwords derived from larger morphological families elicited larger N250 amplitudes than those from smaller families, $\beta = 0.089$, $SE = 0.025$, $t(4851) = 3.53$, $p < .001$. Critically, early morphological processing was strongly modulated by individual differences. Morphological Complexity interacted robustly with Orthographic Sensitivity, $\beta = -0.196$, $SE = 0.028$, $t(4851) = -6.98$, $p < .001$, and with Language Proficiency, $\beta = -0.179$, $SE = 0.020$, $t(4851) = -8.76$, $p < .001$, indicating that the magnitude of the complexity effect at the N250 varied substantially as a function of reader skill. In addition, the Complexity × Family Size × Orthographic Sensitivity interaction was significant, $\beta = -0.065$, $SE = 0.028$, $t(4851) = -2.32$, $p = .021$, suggesting that the influence of morphological family structure on early processing of complex nonwords depended on orthographic sensitivity. No corresponding three-way interaction involving Language Proficiency was observed. Together, these results indicate that early morpho-orthographic parsing of novel forms is highly sensitive to individual differences, with both orthographic sensitivity and language proficiency systematically shaping complexity-related processing costs at the N250.

To clarify the significant Morphological Complexity × Orthographic Sensitivity interaction at the N250 for nonwords, we examined estimated marginal means of complexity at three levels of orthographic sensitivity (-1, 0, +1 SD), averaged over family size. For participants with low orthographic sensitivity (-1 SD), complex nonwords elicited substantially more positive N250 amplitudes than simple nonwords, $\Delta = 0.64,\mu V$, $SE = 0.08$, $z = 8.34$, $p < .0001$. A similar but attenuated complexity effect was observed at mean orthographic sensitivity (0 SD), $\Delta = 0.25,\mu V$, $SE = 0.05$, $z = 4.98$, $p < .0001$. In contrast, for participants with high orthographic sensitivity (+1 SD), the complexity effect was reversed in direction and did not reach significance, $\Delta = -0.14,\mu V$, $SE = 0.07$, $z = -1.93$, $p = .053$. Thus, increasing orthographic sensitivity was associated with a progressive reduction—and eventual elimination—of the early processing cost associated with morphological complexity. This pattern indicates that early morpho-orthographic parsing of novel forms is most effortful for readers with lower orthographic sensitivity and becomes increasingly efficient as orthographic sensitivity increases.

We conducted parallel follow-up analyses to characterize the significant Morphological Complexity × Language Proficiency interaction at the N250, again examining estimated marginal means of complexity at -1, 0, and +1 SD of language proficiency, averaged over family size. For participants with low language proficiency (-1 SD), complex nonwords elicited markedly more positive N250 amplitudes than simple nonwords, $\Delta = 0.57,\mu V$, $SE = 0.06$, $z = 9.05$, $p < .0001$. This complexity effect was reduced but remained reliable at mean language proficiency (0 SD), $\Delta = 0.22,\mu V$, $SE = 0.05$, $z = 4.31$, $p < .0001$. At high language proficiency (+1 SD), the complexity effect reversed in direction and reached significance, with simple nonwords eliciting more positive N250 amplitudes than complex nonwords, $\Delta = -0.14,\mu V$, $SE = 0.07$, $z = -2.15$, $p = .032$. These results indicate that increasing language proficiency is associated with a qualitative shift in early processing of morphologically complex nonwords, from enhanced complexity-related costs in lower-proficiency readers to attenuated or reversed effects in higher-proficiency readers.

\paragraph*{Nonwords}
Mean N400 amplitudes for nonword stimuli were analyzed using a linear mixed-effects model with fixed effects of Morphological Complexity, Family Size, Orthographic Sensitivity (OS), and Language Proficiency (LP), including all interactions among these factors. Random intercepts were included for participants and for electrode site nested within participant. The analysis revealed a robust main effect of Morphological Complexity, such that complex nonwords elicited more positive (less negative) N400 amplitudes than simple nonwords, $\beta = 0.235$, $SE = 0.031$, $t(4851) = 7.64$, $p < .001$. There was no reliable main effect of Family Size, $\beta = 0.042$, $SE = 0.031$, $t(4851) = 1.38$, $p = .17$, nor of Language Proficiency or Orthographic Sensitivity (both $p > .28$). Critically, Morphological Complexity interacted with both individual-difference measures. The Complexity × Language Proficiency interaction was significant, $\beta = -0.127$, $SE = 0.025$, $t(4851) = -5.08$, $p < .001$, as was the Complexity × Orthographic Sensitivity interaction, $\beta = -0.120$, $SE = 0.034$, $t(4851) = -3.49$, $p < .001$, indicating that the magnitude and direction of N400 complexity effects varied systematically as a function of reader skill. In addition, Morphological Complexity interacted with Family Size, $\beta = 0.111$, $SE = 0.031$, $t(4851) = 3.63$, $p < .001$. This interaction was further modulated by Language Proficiency, yielding a significant Complexity × Family Size × Language Proficiency interaction, $\beta = -0.057$, $SE = 0.025$, $t(4851) = -2.29$, $p = .022$. The corresponding three-way interaction involving Orthographic Sensitivity was not significant, $p = .30$.

Follow-up analyses examined estimated marginal means to characterize the Complexity × Language Proficiency interaction. For readers with low language proficiency (-1 SD), complex nonwords elicited substantially more positive N400 amplitudes than simple nonwords, particularly for items from large morphological families, consistent with strong semantic co-activation. At mean proficiency, this complexity effect was attenuated, and at high language proficiency (+1 SD), the complexity effect was markedly reduced or absent across family sizes. Direct contrasts confirmed that the joint influence of complexity and family size was strongest for low-proficiency readers (Complex - Simple × Large - Small: $\Delta = 0.67,\mu V$, $SE = 0.16$, $z = 4.31$, $p < .0001$) and substantially weaker and non-significant for high-proficiency readers ($\Delta = 0.21,\mu V$, $SE = 0.16$, $z = 1.31$, $p = .19$).

Parallel follow-up analyses for Orthographic Sensitivity revealed a similar but less pronounced pattern. Readers with low orthographic sensitivity showed large N400 complexity effects, particularly for nonwords derived from large families, whereas these effects were attenuated at mean sensitivity and minimal at high sensitivity. This contrast between the two individual difference dimensions at the N400 is theoretically significant. Language Proficiency specifically modulated the joint influence of morphological complexity and family size — that is, the degree to which semantic co-activation from morphological relatives shaped processing of novel forms — which is precisely the pattern predicted if Language Proficiency indexes the depth of morphological family representations in semantic space. Orthographic Sensitivity, by contrast, influenced overall sensitivity to complexity without specifically shaping the contribution of family size, suggesting that it affects the general efficiency of semantic access for novel forms rather than the structure of morphological family activation per se.

\subsection{Results Summary}
Across behavioral and electrophysiological measures, the results reveal a consistent dissociation between the sensitivity of behavioral and neural measures to individual differences, alongside a partial dissociation between the roles of Orthographic Sensitivity and Language Proficiency in modulating morphological processing dynamics.

Behaviorally, response times were robustly driven by item-level morphological structure — base frequency and family size for words, morphological complexity and base frequency for nonwords — but neither individual difference dimension selectively modulated the structure of these effects. Orthographic Sensitivity predicted overall response speed for nonwords, and showed a marginal trend for words, suggesting that spelling-based skill contributes to general processing efficiency without altering the specific morphological effects that govern response times.

ERP measures told a different story. At the N250, family size modulated early processing for words, an effect not straightforwardly predicted by multi-stage accounts. For nonwords, both Orthographic Sensitivity and Language Proficiency modulated the complexity effect at the N250, indicating that early morphological parsing of novel forms is sensitive to accumulated lexical experience across both dimensions rather than specifically to orthographic skill as predicted. At the N400, the predicted role of Language Proficiency in later semantic integration received its clearest support: Language Proficiency specifically modulated the joint influence of morphological complexity and family size for nonwords, consistent with the hypothesis that deeper lexical-semantic knowledge shapes the degree to which morphological family structure is recruited during semantic integration. Orthographic Sensitivity affected overall N400 sensitivity to complexity for nonwords but did not specifically modulate the family size contribution, suggesting a more general influence on semantic access efficiency rather than on the structure of morphological family activation.

Two findings deserve particular emphasis. First, the dissociation between behavioral and ERP sensitivity to individual differences indicates that reader skill modulates the neural dynamics of morphological processing even when behavioral responses appear relatively uniform — a finding that underscores the importance of neural measures for characterising individual differences that are invisible at the level of response times. Second, while the predicted clean dissociation between Orthographic Sensitivity and Language Proficiency was only partially supported — both dimensions influencing early processing of novel forms — the specific role of Language Proficiency in modulating the family size contribution to semantic integration at the N400 represents the clearest evidence to date that depth of lexical-semantic knowledge shapes a qualitatively distinct aspect of morphological processing. The theoretical implications of this pattern are developed in the Discussion.


\section{Discussion}

\subsection{Overview of Findings}
The present study examined how morphological structure and item-level lexical familiarity shape visual word recognition, and how these effects are modulated by individual differences in reading-related skill. By combining behavioral measures with time-resolved electrophysiological indices, we were able to distinguish effects that reflect overall processing efficiency from those that reveal how morphological information is constructed and integrated over time. Across analyses, a clear dissociation emerged: reaction times were driven primarily by item-level lexical familiarity and morphological complexity, whereas ERPs revealed robust, stage-specific modulation by Orthographic Sensitivity and Language Proficiency, particularly for novel forms. These findings suggest that individual differences exert their strongest influence not on whether morphology is engaged, but on how and when morphological structure is parsed and regulated during lexical processing.

\subsection{Behavioral Evidence: Efficiency Without Structural Modulation}
Behavioral response times showed robust and theoretically coherent effects of item-level lexical familiarity (Base Frequency and Family Size) and morphological complexity. For words, higher base frequency and larger family size independently facilitated lexical decisions, consistent with prior evidence that stem familiarity and morphological connectivity contribute additively to recognition speed. For nonwords, morphological complexity reliably slowed responses, and higher base frequency produced further interference, indicating that morphological parsing is at least partially automatic and can hinder decision-making when lexical legitimacy is absent.

Critically, however, individual differences did not selectively modulate these structural effects. Orthographic Sensitivity predicted faster responses overall, reflecting greater efficiency in decoding and form-based processing, but did not alter the magnitude of morphological effects. Language Proficiency likewise showed little impact on response speed. Together, these findings indicate that behavioral measures primarily index processing efficiency and familiarity, rather than differences in the structure or timing of morphological computation across readers. This absence of behavioral modulation contrasts sharply with the ERP results and underscores the importance of neural measures for revealing latent differences in processing dynamics.

ubsection{Neural Evidence for Stage-Specific Modulation by Individual Differences}

\subsubsection{Early Parsing (N250): Skill-Dependent Construction of Morphological Structure}
The N250 component, associated with early morpho-orthographic parsing, revealed a striking divergence between words and nonwords. For words, N250 amplitudes were reliably modulated by Family Size, indicating early access to stored morphological representations. Importantly, this effect showed little sensitivity to individual differences, suggesting that when morphological structure is lexically entrenched, early parsing proceeds in a relatively uniform manner across readers.

In contrast, N250 responses to nonwords were strongly shaped by both Orthographic Sensitivity and Language Proficiency. Morphological complexity elicited large early processing costs in lower-skill readers, whereas these effects were attenuated or reversed in higher-skill readers. This pattern suggests that when morphological structure must be constructed online—as is the case for novel forms—early parsing is highly sensitive to reader skill. More skilled readers appear able to segment and evaluate morphological structure efficiently, minimizing early processing costs, whereas less skilled readers engage in broader or less selective decomposition that amplifies complexity-related effects.

\subsubsection{Later Integration (N400): Proficiency-Dependent Semantic Regulation}
At the N400, which indexes semantic integration, individual differences exerted even more differentiated effects. For words, N400 amplitudes were primarily shaped by Family Size, consistent with facilitated integration for items embedded in richer morphological networks. Orthographic Sensitivity exerted an independent main effect, but neither individual-difference measure reliably modulated family-size effects, indicating that semantic integration of familiar words is relatively stable across readers.

For nonwords, however, Language Proficiency played a central role. Morphological complexity interacted robustly with Language Proficiency, and this interaction was further qualified by Family Size. Low-proficiency readers exhibited strong N400 complexity effects, particularly for nonwords derived from large families, suggesting heightened semantic co-activation that may be difficult to regulate when meaning fails to materialize. In contrast, high-proficiency readers showed attenuated or absent complexity effects, indicating more effective suppression of misleading morphological activation. Orthographic Sensitivity showed a parallel but weaker pattern, influencing overall sensitivity to complexity without reliably modulating the joint influence of complexity and family structure.

Together, these findings indicate that Language Proficiency is particularly critical for regulating semantic activation during later stages of processing, especially when morphological structure is suggestive but semantically unlicensed.

\subsection{Dissociable Roles of Orthographic Sensitivity and Language Proficiency}

The present results clarify the complementary but distinct contributions of Orthographic Sensitivity and Language Proficiency to morphological processing. Orthographic Sensitivity primarily reflects efficiency in decoding and form-based analysis, conferring global speed advantages in behavior and attenuating early complexity costs in the N250. Its influence is strongest when readers must rapidly parse unfamiliar strings into plausible sub-lexical units.

Language Proficiency, by contrast, exerts minimal influence on overt response speed but profoundly shapes neural responses, particularly at later stages of processing. High-proficiency readers appear better able to regulate semantic activation, enhancing integration when morphological structure aligns with meaning and suppressing it when structure is misleading. This dissociation explains why Language Proficiency yields robust ERP effects in the absence of corresponding behavioral modulation: lexical decision tasks do not require deep semantic resolution, but neural measures reveal differences in how meaning is evaluated and controlled.

\subsection{Implications for Models of Morphological Processing}

These findings support an interactive, multi-stage view of morphological processing in which decomposition is obligatory but its consequences are dynamically regulated. Item-level lexical familiarity drives efficient access and integration, producing additive behavioral effects, whereas individual differences tune the balance between automatic parsing and controlled semantic evaluation. The strong modulation of nonword ERPs by reader skill highlights the importance of considering individual variability when modeling morphological processing, particularly in contexts where structure must be inferred rather than retrieved.

Rather than reflecting fixed processing strategies, morphological effects emerge from the interaction between stimulus structure and reader skill, with early parsing and later integration differentially sensitive to orthographic and linguistic experience.

\subsection{Limitations and Future Directions}

The present study focused on visual word recognition in skilled adult readers; future work could extend these findings by examining developmental populations or by manipulating semantic transparency to further dissociate form-based and meaning-based contributions. Longitudinal or training studies may also clarify how orthographic sensitivity and language proficiency jointly evolve and reshape the temporal dynamics of morphological processing

\subsection{Conclusion}

In sum, the present findings demonstrate that individual differences in reading skill do not primarily alter whether morphology is engaged, but rather shape how morphological structure is constructed and regulated over time. Behavioral measures reflect efficiency and familiarity, whereas neural measures reveal dynamic, skill-dependent modulation of early parsing and later integration. These results underscore the value of combining behavioral and electrophysiological approaches to capture the adaptive nature of morphological processing in skilled reading.

\subsection{Figures}
Frontiers requires figures to be submitted individually, in the same order as they are referred to in the manuscript. Figures will then be automatically embedded at the bottom of the submitted manuscript. Kindly ensure that each table and figure is mentioned in the text and in numerical order. Figures must be of sufficient resolution for publication \href{https://www.frontiersin.org/about/author-guidelines#ImageSizeRequirements}{see here for examples and minimum requirements}. Figures which are not according to the guidelines will cause substantial delay during the production process. Please see \href{https://www.frontiersin.org/about/author-guidelines#FigureRequirementsStyleGuidelines}{here} for full figure guidelines. Cite figures with subfigures as figure \ref{fig:Subfigure 1} and \ref{fig:Subfigure 2}.


\subsection{Tables}
Tables should be inserted at the end of the manuscript. Please build your table directly in LaTeX.Tables provided as jpeg/tiff files will not be accepted. Please note that very large tables (covering several pages) cannot be included in the final PDF for reasons of space. These tables will be published as \href{http://home.frontiersin.org/about/author-guidelines#SupplementaryMaterial}{Supplementary Material} on the online article page at the time of acceptance. The author will be notified during the typesetting of the final article if this is the case. 


\section{Additional Requirements}

For additional requirements for specific article types and further information please refer to \href{http://www.frontiersin.org/about/AuthorGuidelines#AdditionalRequirements}{Author Guidelines}.

\section*{Conflict of Interest Statement}
The authors declare that the research was conducted in the absence of any commercial or financial relationships that could be construed as a potential conflict of interest.

\section*{Author Contributions}
JM: Conceptualization, Methodology, Software, Validation, Formal analysis, Resources, Data curation, Visualization, Writing – original draft, Writing – review \& editing, Supervision.
EK: Methodology, Software, Investigation, Data curation, Writing – original draft, Writing – review \& editing, Supervision.
FG: Methodology, Software, Investigation, Data curation.
JP: Investigation, Data curation.
TR: Investigation, Data curation.
EM: Investigation, Data curation.
CR-M: Investigation, Data curation.
NA: Investigation, Data curation.

All authors contributed to the article and approved the submitted version.

\section*{Funding}
Research reported in this manuscript was supported by the Rhode Island IDeA Network of Biomedical Research Excellence, funded by the National Institute of General Medical Sciences of the National Institutes of Health under grant number P20GM103430. Additional support for student researchers (Joemari Pulido, Tess Rooney, Natalia Alzate, Ethan Moore, and Crismar Ramos-Marte) and for data collection was provided through the Center for Engaged Learning at Providence College and a series of Hampshire College Faculty Development Grants awarded to Dr. Morris during her tenure at that institution. The funders had no role in study design, data collection and analysis, decision to publish, or preparation of the manuscript.


\section*{Acknowledgments}
We thank Sophie Fulghum, Dylan Palumbo, Maeve Garrity, Aleksandra Kurylowicz, Sarah Debian, Aleksandra Saneff, Kenny Diaz, and Kiley Vallee for their assistance with participant recruitment and data collection.

\section*{Supplemental Data}
 \href{http://home.frontiersin.org/about/author-guidelines#SupplementaryMaterial}{Supplementary Material} should be uploaded separately on submission, if there are Supplementary Figures, please include the caption in the same file as the figure. LaTeX Supplementary Material templates can be found in the Frontiers LaTeX folder.

\section*{Data Availability Statement}
The datasets [GENERATED/ANALYZED] for this study can be found in the [NAME OF REPOSITORY] [LINK].
% Please see the availability of data guidelines for more information, at https://www.frontiersin.org/about/author-guidelines#AvailabilityofData

\bibliographystyle{Frontiers-Harvard} %  Many Frontiers journals use the Harvard referencing system (Author-date), to find the style and resources for the journal you are submitting to: https://zendesk.frontiersin.org/hc/en-us/articles/360017860337-Frontiers-Reference-Styles-by-Journal. For Humanities and Social Sciences articles please include page numbers in the in-text citations 
%\bibliographystyle{Frontiers-Vancouver} % Many Frontiers journals use the numbered referencing system, to find the style and resources for the journal you are submitting to: https://zendesk.frontiersin.org/hc/en-us/articles/360017860337-Frontiers-Reference-Styles-by-Journal
\bibliography{1_VSL_INDIFFS.bib}


\end{document}
